
%% ===================================================================
%% LP-134 canonical naming (2025-12-15):
%% Per LP-134 (Lux Chain Topology), M-Chain has been split into:
%%   M-Chain — MPC ceremonies for bridge custody of EXTERNAL wallets
%%             (CGGMP21, FROST, Pulsar-general)
%%   F-Chain — FHE compute + TFHE bootstrap-key generation (encrypted EVM)
%% M-Chain is REMOVED entirely; teleport is B-Chain (bridgevm, LP-6000). No teleportvm.
%% Where this paper says "M-Chain MPC" / "M-Chain threshold governance" /
%% "M-Chain TFHE" / "M-Chain custody", read as M-Chain (MPC) or F-Chain (FHE).
%% ===================================================================

\section{Security Model}
\label{sec:security}

The security model addresses three threat classes: agent compromise, principal compromise, and infrastructure compromise.

\subsection{Agent Compromise Containment}

\begin{theorem}[Capsule Isolation]
A compromised Agent Capsule $\mathcal{C}_A$ cannot affect the state, budget, or capabilities of any other capsule $\mathcal{C}_B$ where $A \neq B$, provided $A$ has not been granted delegation authority over $B$'s resources.
\end{theorem}

\begin{proof}
Each capsule's vault is encrypted with a key derived from its own DID private key (Equation~\ref{eq:vault-key}). Cross-capsule access requires a valid capability token signed by the target capsule's principal. A compromised capsule can forge signatures only under its own DID key, which is not accepted by other capsules' verification logic. Budget draws are scoped to the capsule's own budget account by the C-Chain contract. Therefore, compromise of $\mathcal{C}_A$ is contained to $\mathcal{C}_A$'s own resources.
\end{proof}

\subsection{Capability Revocation Under Compromise}

When compromise is detected, the principal revokes the capsule's DID by publishing a revocation entry on the I-Chain. This has three immediate effects:

\begin{enumerate}
    \item All capability tokens referencing the revoked DID become invalid---verifiers reject them upon checking the revocation registry.
    \item The budget account is frozen by the C-Chain contract, which checks DID status before processing withdrawals.
    \item All delegation chains rooted at the compromised capsule are invalidated (Section~\ref{sec:identity}).
\end{enumerate}

\subsection{Threshold Recovery}

If the agent's private key is lost (not compromised), the principal can initiate threshold recovery via the M-Chain. The capsule's vault encryption key was escrowed at provisioning time using $(k, n)$-threshold secret sharing among a recovery committee selected by the principal:

\begin{equation}\label{eq:recovery}
    k_{\text{vault}} = \text{TSS.Reconstruct}\big(\{s_i\}_{i \in S}, |S| \geq k_{\text{rec}}\big)
\end{equation}

Recovery creates a new DID for the capsule, migrates the vault to a new encryption key, and invalidates the old DID. The action history is preserved.

\subsection{Infrastructure Compromise}

The execution environment (cloud VM, container, edge device) is untrusted. The vault is encrypted at rest and in transit. The agent's private key is stored in a hardware security module (HSM) or TEE where available. Even if the host is fully compromised, the attacker obtains only ciphertext and encrypted memory---useless without the DID private key or the threshold recovery committee.

